% 第17批：自测题6填空6.2.1—6.2.9；印刷165页／PDF173页。
\subsection{自测题6：静电场、磁场与麦克斯韦方程填空}

\sourcetest{6.2.1}{165}{173}
\problemtext 均匀带电直线段长 $d$、电荷线密度 $+\lambda$，以线段中点为球心作半径 $R>d$ 的球面，求通量和线段延长线上球面交点 $P$ 的场强。
\workings 高斯定律给出封闭球面的电通量
\[\boxed{\Phi_E=\frac{\lambda d}{\varepsilon_0}}.\]
以线段中心为 $x=0$，$P=(-R,0)$，由沿 $x$ 的库仑积分，
\[\boxed{|\bm E(P)|=\frac{\lambda}{4\pi\varepsilon_0}
\left(\frac1{R-d/2}-\frac1{R+d/2}\right)
=\frac{\lambda d}{4\pi\varepsilon_0(R^2-d^2/4)}}.\]
正电荷时 $\bm E(P)$ 沿负 $x$ 方向。
\crosscheck 球面通量只取决于包围总电荷，但局部场强不能把线段直接替换成中心点电荷；要求 $R>d/2$，题设 $R>d$ 足够。

\sourcetest{6.2.2}{165}{173}
\problemtext 两无限大平行带电面 $A,B$，两面之间场强 $E_0$ 向左；左外侧 $E_0/3$ 向左，右外侧 $E_0/3$ 向右。求两面面电荷密度。
\workings 令右向为正。跨越电荷面场法向跳跃满足 $E_{\rm right}-E_{\rm left}=\sigma/\varepsilon_0$。面 $A$：$\sigma_A=\varepsilon_0(-E_0+E_0/3)$；面 $B$：$\sigma_B=\varepsilon_0(E_0/3+E_0)$，故
\[\boxed{\sigma_A=-\frac23\varepsilon_0E_0,\qquad \sigma_B=+\frac43\varepsilon_0E_0}.\]
\crosscheck 带电片不是必然等量异号；外侧场非零意味着总面电荷密度和不为零。

\sourcetest{6.2.3}{165}{173}
\problemtext 负电荷由无穷远靠近一块不带净电荷、与地绝缘的导体，问导体内部电场以及导体电势怎样变化。
\workings 静电平衡时导体内部 $\boxed{\bm E=0}$，故场强保持零。外部负电荷的靠近，使原本零参考电势的孤立中性导体整体感应到较低电势，故 \textbf{电势减小}。
\crosscheck 导体即使净电荷保持零，也会发生表面电荷重新分布；“孤立”与“接地”边界条件不同。

\sourcetest{6.2.4}{165}{173}
\problemtext 两平行无限长直导线 $A,B$ 电流大小相同、方向均垂直纸面向外、间距 $a$，$P$ 为二者中点。另有图示仅包围 $A$ 的闭合回路 $l$，求 $P$ 处磁感应强度及闭路积分。
\workings 在中点，两条直线对 $P$ 的磁场方向相反、模相等，故 $\boxed{\bm B(P)=0}$。安培环路定律以面法向朝纸面外为正，图中回路只围住 $A$：
\[\boxed{\oint_l\bm B\cdot d\bm l=\mu_0 I}\]
（若取反向绕行则为负）。
\crosscheck 不能由一点 $\bm B(P)=0$ 断言整条回路积分为零；必须数穿过回路围成面的总电流。

\sourcetest{6.2.5}{165}{173}
\problemtext 平均周长0.50m 的闭合细铁环绕有500匝线圈，电流0.30A，铁芯相对磁导率600，求环中 $H,B$。
\workings 环路定律给出 $\oint\bm H\cdot d\bm l=NI$；在细环近似下
\[\boxed{H=\frac{NI}{\ell}=\frac{500(0.3)}{0.50}=300\ {\rm A/m}}.\]
且 $B=\mu_0\mu_rH$，故
\[\boxed{B\simeq(4\pi\times10^{-7})(600)(300)=0.226\ {\rm T}}.\]
\crosscheck 铁磁性材料的相对磁导率可能非线性变化，此题按题给常数 $\mu_r$ 的线性近似。

\sourcetest{6.2.6}{165}{173}
\problemtext 条形磁铁插入与冲击电流计串联的金属环，经过仪器的总感应电荷量 $q=2.0\times10^{-5}$C，回路总电阻25Ω，求磁通变化量。
\workings 法拉第定律 $\mathcal E=-d\Phi/dt$，$I=\mathcal E/R$。积分得 $q=\int I\,dt=-\Delta\Phi/R$，所以
\[\boxed{|\Delta\Phi|=Rq=25(2.0\times10^{-5})=5.0\times10^{-4}\ {\rm Wb}}.\]
\crosscheck 电荷量通常给的是累计流过电量的正值，因此结果报告磁通变化量的\emph{大小}；有向符号由线圈取向定义。

\sourcetest{6.2.7}{165}{173}
\problemtext 一根向上通电的无限长直导线右侧，水平金属棒 $AB$ 与导线共面，且沿导线方向向上以速度 $v$ 移动。比较两端感应电势，并讨论电流反向或棒改为平行导线的情况。
\workings 原电流向上时，右侧磁场指入纸面，$\bm v\times\bm B$ 向左，故 \textbf{$V_A>V_B$}；电流反向则 \textbf{$V_B>V_A$}。若棒沿导线方向放置，$\bm v\times\bm B$ 仍水平且与棒正交，因此 $\boxed{V_A=V_B}$（仅指运动电动势贡献）。
\crosscheck 这里比较的是运动棒中电荷受洛伦兹磁力分离产生的电势差；同时假定导体无外加回路电压。

\sourcetest{6.2.8}{165}{173}
\problemtext 无限长直导线载电流 $I$，求距轴线 $a$ 处真空磁场的能量密度。
\workings $B(a)=\mu_0I/(2\pi a)$，真空磁场能量密度
\[\boxed{w_m=\frac{B^2}{2\mu_0}=\frac{\mu_0I^2}{8\pi^2a^2}}.\]
\crosscheck 该式是局部能量密度；若求整条无限长导线外部总能量，还需要外半径截断，否则发散。

\sourcetest{6.2.9}{165}{173}
\problemtext 给定麦克斯韦积分式①电场高斯定律、②法拉第定律、③磁场高斯定律、④安培—麦克斯韦定律；判断“变化的磁场伴随电场”“磁感线没有起止”“电荷伴随电场”各对应哪式。
\workings 由 $\oint\bm E\cdot d\bm l=-d\Phi_B/dt$ 得变化磁场伴随涡旋电场，选 \textbf{②}；$\oint\bm B\cdot d\bm S=0$ 表明磁感线无始无终，选 \textbf{③}；$\oint\bm D\cdot d\bm S=Q_{\rm enclosed}$ 表明电荷是电场散度源，选 \textbf{①}。
\[\boxed{(1)\ \text{②};\quad(2)\ \text{③};\quad(3)\ \text{①}}\]
\crosscheck 方程①表示有电荷必有电场结构，并不意味电荷附近\emph{每一点}的电场都非零；安培—麦克斯韦定律④反映电流和变化电场产生磁场。
